\stopcontents[chapters]
\chapter*{Conclusion}
\label{chap:conclu}

\epigraph{\itshape The easiest way to stop piracy is not by putting antipiracy technology to work. It's by giving those people a service that's better than what they're receiving from the pirates.}{-- Gabe NEWELL}

\addcontentsline{toc}{chapter}{Conclusion}
\fancyhead{}
\fancyhead[LE,RO]{Conclusion}

The investigation carried out across the three preceding chapters converges on a single, structurally grounded answer: cheating in a competitive FPS can be addressed effectively without requiring the privileged, system-wide access that kernel-level anti-cheat currently demands. Chapter 1 established that the sophistication of modern cheats — from input-layer automation to hardware-assisted DMA vectors — has progressively driven anti-cheat design toward Ring 0, and Chapter 2 showed that this escalation carries costs disproportionate to what it protects: an asymmetric arms race that structurally favours the attacker [1], a categorical blind spot against hardware-level manipulation, and a level of local intrusiveness — demonstrated firsthand through the Windows-Kernel-Explorer experiment rather than argued abstractly — that formal analysis has repeatedly likened to rootkit behaviour [3]. Chapter 3 tested the alternative this critique implies. The gameplay analytics pipeline developed there does not merely propose server-side detection as a theoretical possibility; it demonstrates, on a calibrated synthetic population, that a purely telemetry-based system can separate cheaters from legitimate players with zero false positives, provided detection is reframed as a ranked prioritisation problem for human review rather than an autonomous verdict. The perfect AUC-ROC obtained, together with a recall deliberately capped by a conservative threshold, illustrates concretely what this reframing buys: every cheater in the dataset remains correctly ranked above every legitimate player, even where the fixed cut leaves half of them unflagged, so nothing is lost that a review queue could not still recover.

The consistency-based design at the core of this pipeline also surfaces a dynamic with no equivalent on the kernel-level side. Because the CONSISTENCY\_BONUS mechanism only responds to mechanical anomalies that persist across a large share of a player's match history, a cheater seeking to stay under its radar is pushed toward deliberately throttling their tool's aggressiveness — capping aimbot precision, disabling wallhack overlays intermittently, limiting reaction-time assistance to a minority of engagements. Carried far enough, this evasion strategy converges the cheater's observable performance toward the very skilled-player distribution the system is calibrated against, at which point the residual competitive advantage gained from cheating narrows accordingly. This is a materially different asymmetry from the one documented in Chapter 2: where kernel-level evasion (BYOVD, DKOM) leaves the cheat's effectiveness fully intact while only its detectability is contested, evading a behavioural, consistency-based detector degrades the cheat's usefulness as a direct condition of evading it. The incentive structure this creates — rather than any single detection event — is arguably the more durable defence the two paradigms are compared on in Section 3.5.

This conclusion should nonetheless be read against the limits established throughout Chapter 3. The results reported rest on a synthetic dataset whose profile distributions, however carefully calibrated against community benchmarks and prior work [2, 4], cannot reproduce the full behavioural irregularity of authentic play, and the LSTM component evaluated here remains an architectural demonstration rather than a trained classifier. Adapting the pipeline to genuinely novel cheating methods — cheat software specifically engineered to mimic the consistency profile of legitimate skilled play, for instance — was not tested, and would require the very real-match data this thesis was unable to obtain.

None of this weakens the central architectural claim, however: server-side behavioural analysis does not need to outperform kernel-level anti-cheat on every axis to justify its place in a defence stack. Positioned as one layer among several — lightweight client-side checks for the host-resident manipulation kernel drivers already handle well, behavioural scoring for what they cannot reach, and human review as the final arbiter rather than an automated ban — it offers a materially less intrusive path to the same competitive integrity objective, one that does not require every player to surrender Ring 0 visibility over their machine simply to be allowed to play.
